\documentclass[10pt,a4paper]{amsart}
\usepackage[margin=19mm]{geometry}
\usepackage{amsmath,amssymb,mathtools}
\usepackage[scr=rsfs,cal=euler]{mathalfa}
\usepackage{xfrac}
\usepackage[unicode,hidelinks,bookmarks=false]{hyperref}
\newcommand{\ie}{i.e.\ }
\newcommand{\cf}{cf.\ }
\newcommand{\resp}{respectively\ }
\newcommand{\Subsection}{\subsection}
\providecommand{\nobreakdash}{\nobreak\hbox{-}}
\setlength{\emergencystretch}{2em}
\multlinegap0pt
\usepackage{amssymb}
\newtheorem{lemma}{Lemma}
\title{Existence of stable periodic orbits in billiards close to lemon and moon billiards}
\author{Alexander Grigo}
\date{Source: arXiv:2603.00275v1 (CC BY 4.0)}
\begin{document}
\maketitle

\noindent\emph{Source and licence.} Existence of stable periodic orbits in billiards close to lemon and moon billiards. Version arXiv:2603.00275v1 (CC BY 4.0), distributed by arXiv under the Creative Commons Attribution 4.0 International licence, http://creativecommons.org/licenses/by/4.0/, as recorded in the arXiv licence metadata for this exact version and shown on the abstract page https://arxiv.org/abs/2603.00275v1. Full licence notice: \texttt{licenses/arxiv-2603.00275-CC-BY-4.0.txt}.\par\smallskip

\noindent\textit{Editorial scope.} Counted unit: the article's lemma lem\_t0\_2t1\_tc (source0.tex lines 883--900). The article's complete original proof of it is delivered with it (source0.tex lines 901--933, one \texttt{proof} environment of 33 lines). No mathematics is written, repaired or re-derived: the statement and the proof are the article's own lines, in the article's own notation, and no retained line is substituted or re-flowed.  Retained uncounted as the source-local context the proof needs: lines 626--632; lines 698--701; lines 713--718.  Nothing was omitted from any retained span, so the package carries no omission record.  No retained line contains a citation, so the wrapper carries no bibliography and no bibliography entry.  No retained line contains a figure environment, an image-inclusion command or drawing code, so no image is delivered and the package declares no asset dependency.  Editorial formatting only, in addition: an AMS \texttt{amsart} preamble that restates the article's own declarations and macros used by the retained lines, namely the environments \texttt{lemma} on separate counters, together with the standard packages those lines need.  The retained environments print in this document as Lemma 1 in order of appearance, whereas the article prints the same items with the numbering of its own full document.  The complete native source of the article as distributed in the arXiv e-print for this exact version is bundled unchanged as \texttt{sources/195473/source0.tex}.
\begin{equation}
  \label{eqn_phi1_alpha1}
  \alpha_1 = \frac{\pi}{N} + \frac{N-1}{N}\,\varepsilon
  \;,\quad
  \frac{\pi}{2} - \varphi_1 = \frac{\pi}{ N } - \frac{1}{N}\,\varepsilon
  \;.
\end{equation}

\begin{equation}
  \label{eqn_rangle_tau0}
  0 < \tau_0 < 2\,r \sin\alpha_1
\end{equation}

\begin{equation}
  \label{eqn_tau1}
  \tau_1
  =
  \frac{ 2\,r \sin\alpha_1 - \tau_0 }{ 2 \cos\varepsilon }
\end{equation}

\begin{lemma}
  \label{lem_t0_2t1_tc}
  The three free paths $\tau_0$, $\tau_1$, $\tau_c$
  satisfy
  \begin{equation*}
    2\,\tau_1 + \tau_0 - \tau_c
    =
    \Big(\frac{1}{ \cos\varepsilon } - 1 \Big)
    \,[ 2\,r \sin\alpha_1 - \tau_0 ]
    + 2\,r\,[ \sin\alpha_1 - \sin(\alpha_1 - \varepsilon) ]
    >
    0
  \end{equation*}
  for all $\varepsilon>0$ (small) and all admissible $\tau_0$.
  In particular,
  $2\,\tau_1 + \tau_0 - \tau_c = \operatorname{\mathcal{O}}(\varepsilon)$
  holds uniformly in $\varepsilon$.
\end{lemma}

\begin{proof}
  By
  \eqref{eqn_tau1}
  and the definition of $\tau_c$ we have
  \begin{equation*}
    2\,\tau_1 + \tau_0 - \tau_c
    =
    \frac{ 2\,r \sin\alpha_1 - \tau_0 }{ \cos\varepsilon }
    + \tau_0
    - 2\,r\cos\varphi_1
    \;.
  \end{equation*}
  Recall also that by \eqref{eqn_phi1_alpha1}
  $ \frac{\pi}{2} - \varphi_1 = \alpha_1 - \varepsilon $, so that
  $\cos\varphi_1 = \sin(\alpha_1 - \varepsilon)$, and therefore
  \begin{equation*}
    2\,\tau_1 + \tau_0 - \tau_c
    =
    \frac{ 2\,r \sin\alpha_1 - \tau_0 }{ \cos\varepsilon }
    + \tau_0
    - 2\,r \sin(\alpha_1 - \varepsilon)
  \end{equation*}
  holds, which can be readily seen to be equivalent to the claimed
  expression for
  $2\,\tau_1 + \tau_0 - \tau_c$.

  Now, by the restriction \eqref{eqn_rangle_tau0} on the range
  for $\tau_0$ it follows that the first term on the right-hand-side
  of the claimed expression for
  $2\,\tau_1 + \tau_0 - \tau_c$
  is positive. The second term is obviously also positive for
  any $\varepsilon>0$ (small).
\end{proof}

\end{document}
